\documentclass[UTF8,fontset=fandol,11pt]{ctexart}
\usepackage[a4paper,margin=24mm]{geometry}
\usepackage{amsmath,amssymb,amsthm,graphicx,tabularx,booktabs,hyperref}
\newtheorem{definition}{定义}[section]
\newtheorem{theorem}[definition]{定理}
\newtheorem{example}[definition]{例题}
\newtheorem{exercise}[definition]{练习}
\hypersetup{unicode=true,colorlinks=true,linkcolor=blue}
\title{数列极限入门：定义、唯一性与简单运算}
\date{}
\begin{document}
\maketitle
\section{数列极限的定义与唯一性}\label{sec:S01}
% UNIT S01-U01
\begin{definition}
\textbf{数列极限}\par

设 $\{a_n\}$ 是定义在 $\mathbb{N}^+=\{1,2,3,\ldots\}$ 上的实数数列，$A\in\mathbb{R}$。直观地说，当下标 $n$ 足够大时，$a_n$ 可以任意接近 $A$，就称数列 $\{a_n\}$ 收敛于 $A$，记作
\[
\lim_{n\to\infty}a_n=A.
\]
严格地说，若对任意 $\varepsilon>0$，都存在正整数 $N$，使得对所有满足 $n>N$ 的正整数 $n$，均有
\[
|a_n-A|<\varepsilon,
\]
则称数列 $\{a_n\}$ 收敛于 $A$。
\label{def:sequence-limit}


这里的量词顺序是：先任意给定 $\varepsilon>0$，再选择正整数 $N$，最后要求所有 $n>N$ 都满足不等式。所选择的 $N$ 可以依赖于 $\varepsilon$，但不能在检验不同的 $n$ 时随 $n$ 改变。
\end{definition}
% UNIT S01-U02
\begin{theorem}
\textbf{数列极限的唯一性}\par

若数列 $\{a_n\}$ 收敛于 $A$，且也收敛于 $B$，其中 $A,B\in\mathbb{R}$，则
\[
A=B.
\]
\label{thm:limit-uniqueness}

\end{theorem}
% UNIT S01-U03
\begin{proof}

反设 $A\ne B$。取
\[
\varepsilon=\frac{|A-B|}{3}>0.
\]
由 $a_n\to A$，存在正整数 $N_1$，使得当 $n>N_1$ 时，
\[
|a_n-A|<\varepsilon.
\]
同理由 $a_n\to B$，存在正整数 $N_2$，使得当 $n>N_2$ 时，
\[
|a_n-B|<\varepsilon.
\]
取正整数 $n>\max\{N_1,N_2\}$，则由三角不等式，
\[
|A-B|=|(A-a_n)+(a_n-B)|
\le |A-a_n|+|a_n-B|
<2\varepsilon=\frac{2|A-B|}{3}.
\]
这意味着 $|A-B|<\frac{2|A-B|}{3}$，与 $|A-B|>0$ 矛盾。因此 $A=B$。

\end{proof}
% UNIT S01-U04
\begin{example}

用数列极限的定义证明
\[
\lim_{n\to\infty}\frac{1}{n}=0.
\]

\end{example}
% UNIT S01-U05
\par\noindent\textbf{解：}
任意给定 $\varepsilon>0$。为了使
\[
\left|\frac{1}{n}-0\right|<\varepsilon,
\]
只需使 $\frac{1}{n}<\varepsilon$，即 $n>\frac{1}{\varepsilon}$。因此，取正整数 $N$ 满足
\[
N>\frac{1}{\varepsilon}.
\]
例如可以取 $N=\left\lfloor\frac{1}{\varepsilon}\right\rfloor+1$。当任意正整数 $n>N$ 时，有
\[
\left|\frac{1}{n}-0\right|=\frac{1}{n}<\frac{1}{N}<\varepsilon.
\]
所以按照数列极限的定义，
\[
\lim_{n\to\infty}\frac{1}{n}=0.
\tag{1}\label{eq:reciprocal-limit}
\]
这里 $N$ 可以依赖于预先给定的 $\varepsilon$，但一旦选定，就必须对所有 $n>N$ 同时成立。
\par
% UNIT S01-U06
\begin{exercise}

用数列极限的定义证明
\[
\lim_{n\to\infty}\left(2+\frac{3}{n}\right)=2.
\]

\end{exercise}
\section*{练习答案}
% UNIT S01-U07
\par\noindent\textbf{解：}
任意给定 $\varepsilon>0$。有
\[
\left|\left(2+\frac{3}{n}\right)-2\right|=\frac{3}{n}.
\]
为了使上式小于 $\varepsilon$，只需使
\[
\frac{3}{n}<\varepsilon,
\qquad\text{即}\qquad
n>\frac{3}{\varepsilon}.
\]
取正整数 $N$ 满足
\[
N>\frac{3}{\varepsilon}.
\]
当任意正整数 $n>N$ 时，便有
\[
\left|\left(2+\frac{3}{n}\right)-2\right|
=\frac{3}{n}
<\frac{3}{N}
<\varepsilon.
\]
因此，根据数列极限的定义，
\[
\lim_{n\to\infty}\left(2+\frac{3}{n}\right)=2.
\]
\par
\end{document}
